\section{Complexity and Tractable Cases}
\label{sec:complexity}

The \emph{budget decision problem} asks whether a $k$-instability witness
from Definition~\ref{def:witness} exists for the given $(k,\theta)$.
The optimization version additionally minimizes its size within that
budget. We first establish fixed-query hardness through source sharing,
then show that singleton supports admit exact algorithms. For the latter,
one may compute the unbudgeted optimum $\rho_\theta$ and compare it
with $k$: an optimum at most $k$ solves \MIW, while a larger optimum
certifies that no witness exists within the budget.

\begin{theorem}[Fixed-query hardness]
\label{thm:hardness}
For each of \agg{COUNT}, \agg{SUM}, \agg{MAX}, \agg{MIN}, and \agg{AVG},
the budget decision problem is NP-complete in data complexity for the
fixed body
\[
R(x)\Join E(x,y,A)\Join R(y).
\]
This holds with values in $\{0,1\}$, designated immutable tuples, and a
self-join of $R$.
\end{theorem}

\begin{figure*}[t]
\centering
\begin{tikzpicture}[
  font=\small,
  vertex/.style={circle,draw,minimum size=7mm,inner sep=0pt},
  selected/.style={vertex,fill=blue!12,draw=blue!65!black,very thick},
  support/.style={draw,rounded corners,minimum width=40mm,
                  minimum height=8mm,fill=black!2},
  flow/.style={-{Stealth[length=2.2mm]},thick,draw=black!55}]
\node[font=\small\bfseries] at (1.6,3.0) {Vertex Cover instance};
\node[vertex] (a) at (0.2,1.8) {$a$};
\node[selected] (b) at (1.6,0.9) {$b$};
\node[vertex] (c) at (3.0,1.8) {$c$};
\draw[thick] (a)--(b)--(c);
\node[align=center] at (1.6,0.05) {$k=1$, cover $C=\{b\}$};

\draw[flow] (3.6,1.5)--(4.5,1.5);
\node[font=\small\bfseries] at (7.0,3.0) {Database instance};
\node[anchor=north] at (5.5,2.5) {\begin{tabular}{c}
  $R$\\\hline
  $a$\\
  \textcolor{blue!65!black}{$\boldsymbol{b}$}\\
  $c$
\end{tabular}};
\node[anchor=north] at (7.9,2.5) {\begin{tabular}{ccc}
  \multicolumn{3}{c}{$E(x,y,A)$}\\\hline
  $a$ & $b$ & $1$\\
  $b$ & $c$ & $1$
\end{tabular}};
\node[align=center] at (7.0,0.05) {$R$ deletable; $E$ immutable};

\draw[flow] (9.4,1.5)--(10.3,1.5);
\node[font=\small\bfseries] at (12.9,3.0) {Derivation supports};
\node[support] at (12.9,2.1)
  {$T_{ab}=\{R(a),\textcolor{blue!65!black}{R(b)}\}$};
\node[support] at (12.9,1.15)
  {$T_{bc}=\{\textcolor{blue!65!black}{R(b)},R(c)\}$};
\node[align=center] at (12.9,0.05)
  {Delete $S=\{R(b)\}$; both die};

\node[draw=black!25,fill=black!3,rounded corners,minimum width=15cm,
      align=center,inner sep=6pt] at (7.55,-0.95)
  {For COUNT and SUM, $Q(D)=2$ and $\theta=2$.\\
   $C$ covers every edge $\Longleftrightarrow$ $S$ destroys every edge derivation;
   $|C|=|S|\le k$.};
\end{tikzpicture}
\caption{A budget-preserving NP-hardness reduction from Vertex Cover for COUNT and SUM,
illustrated with two edges and $k=1$. For a general graph, each vertex
becomes a deletable $R$ tuple and each oriented edge becomes one immutable
$E$ tuple. The fixed body $R(x)\Join E(x,y,A)\Join R(y)$ gives one
support per edge. Setting $\theta=|E_G|$ forces all edge supports to be
hit. The proof extends the construction to MAX, MIN, and AVG using a
permanently surviving derivation; the same budget $k$ is retained.}
\label{fig:vertex-cover}
\Description{Three panels show the reduction. In the left panel, the graph
has edges a-b and b-c, and vertex b alone covers both. The middle panel
contains deletable R tuples a, b, c and immutable E tuples (a,b,1) and
(b,c,1). The right panel shows the supports R(a),R(b) and R(b),R(c).
Deleting the single tuple R(b) destroys both derivations, preserving the
Vertex Cover budget of one.}
\end{figure*}

\begin{proof}
Reduce from Vertex Cover on a simple undirected graph $G=(V,E_G)$ with
at least one edge and Vertex Cover budget $k$. We use exactly the same
budget $k$ in the \MIW\ instance. Orient each edge once. Create a deletable tuple $R(v)$
for each vertex and an immutable tuple $E(u,v,1)$ for each oriented edge.
There is exactly one derivation per edge, with support
$\{R(u),R(v)\}$. A deletion set destroys every such derivation exactly
when its vertices form a vertex cover. Figure~\ref{fig:vertex-cover}
illustrates this correspondence.

For \agg{COUNT} and \agg{SUM}(A), choose $\theta=|E_G|$. The target
is zero. A vertex cover of size at most $k$ therefore exists if and only
if the constructed instance has a $k$-instability witness. For \agg{MAX} and
\agg{AVG}, add fresh vertices $p,q$, immutable $R(p),R(q)$, and
immutable $E(p,q,0)$. This creates one permanently surviving zero-valued
derivation and no additional joins. For \agg{MAX}, use $\theta=1$.
For \agg{AVG}, use $\theta=|E_G|/(|E_G|+1)$, the initial average.
In either case, the target is zero and is reached precisely when every
graph-edge derivation is destroyed. For \agg{MIN}, instead give all
graph edges value zero and the permanent edge value one; choose
$\theta=1$. Reaching minimum one again requires destroying every graph
edge. In every construction, feasible deletion sets and vertex covers
have the same cardinalities, so the budget $k$ is preserved for every
aggregate. The encoding has polynomial size.

For membership in NP, a certificate lists at most $k$ deletable tuples.
Evaluation of this fixed query, the aggregate, and the validity condition
uses polynomial time with exact rational arithmetic. Hence the decision
problem is NP-complete for each operator.
\end{proof}

Only $R$ contains deletable tuples in this construction. Reusing that one
relation in two atom positions is sufficient for hardness across all five
operators. Each support has size two: pairwise coupling inside a single
deletable relation is already sufficient for this reduction.

\begin{proposition}[Singleton-support coverage]
\label{prop:modular}
For \agg{COUNT} or nonnegative \agg{SUM}, suppose $|T_b|\le1$ for
every $b$. Sort deletable tuples by their initial marginal gains. The
shortest prefix reaching the threshold is a minimum $k$-instability
witness if its size is at most $k$; otherwise no witness fits the budget.
Equivalently, feasibility is decided by the total gain of the largest
$\min\{k,n\}$ gains.
\end{proposition}
\begin{proof}
Put $a_t=\sum_{b:T_b=\{t\}}w_b$. The benefit is modular:
$f(S)=\sum_{t\in S}a_t$. Among sets of any fixed cardinality $j$, the
$j$ largest $a_t$ attain the greatest possible benefit. The first feasible
prefix therefore has minimum cardinality. If that prefix exceeds $k$,
no set of at most $k$ deletions can reach the threshold.
\end{proof}
The gains can be computed in $O(M+n)$ operations from explicit singleton
supports, followed by $O(n\log n)$ comparisons for sorting.

\begin{theorem}[Singleton-support tractability]
\label{thm:single-relation}
If $|T_b|\le1$ for every derivation, \MIW\ is solvable in polynomial
time for each of the five aggregates. Consequently, every fixed
self-join-free query with only one deletable relation is tractable.
Allowing that relation to occur twice already admits the fixed-query
NP-complete cases of Theorem~\ref{thm:hardness}.
\end{theorem}
\begin{proof}
Proposition~\ref{prop:modular} handles COUNT and SUM. For an extremum,
each critical singleton forces its tuple's deletion; an empty critical
support makes the instance infeasible. The union of forced tuples is
optimal if any derivation survives, and otherwise no valid witness exists.
If the forced set exceeds $k$, the budgeted instance is infeasible.

For AVG at target $\tau=\tau_{\downarrow}$, define
$F_\tau(S)=\sum_{b\in\B(S)}(v_b-\tau)$ and write
$F_0=F_\tau(\emptyset)>0$.
Let $a_t$ sum $v_b-\tau$ over the derivations with $T_b=\{t\}$.
Singleton supports give
\[
F_\tau(S)=F_0-\sum_{t\in S}a_t.
\]
Every valid threshold-crossing witness preserves a derivation $b_0$ with
$v_{b_0}\le\tau$.
Enumerate its possible support, forbid deleting that support, sort the
remaining positive $a_t$ in descending order, and take the shortest
prefix whose sum reaches $F_0$. This is exact for that protected domain:
nonpositive gains never help, and the largest gains maximize benefit for
each cardinality. Taking the smallest feasible prefix over survivors is
therefore globally optimal before imposing the budget; compare its size
with $k$ to solve the budgeted problem. At most $n+1$ distinct singleton or empty
supports are protected. One initial sort and a scan per protected support
take $O(M+n\log n+n^2)$ operations and comparisons.
\end{proof}
This is a boundary of uniform tractability under a syntactic restriction:
all queries on the self-join-free side are tractable, while the expanded
class contains hard queries. Individual queries with self-joins can still
be tractable. The following independent-row case gives a faster AVG
algorithm when every deletable group has exactly one row.

\begin{proposition}[Independent-row averages]
\label{prop:avg-rows}
For a bag of numeric rows, with each deletable row contributing exactly
one aggregate input independently of other rows, downward \agg{AVG}
\MIW\ is solvable by sorting and scanning prefixes in $O(N+n\log n)$
arithmetic operations and comparisons, where $N$ is the total row count.
\end{proposition}
\begin{proof}
Sort deletable values as $u_1\ge\cdots\ge u_n$. Let $V$ be the initial
sum. For a fixed deletion count $j<N$, the remaining denominator is
$N-j$, so deleting the largest $j$ values minimizes the remaining
average. Check
\[
V-\sum_{i=1}^j u_i\ \le\ \tau_{\downarrow}(N-j)
\]
for each $j\le\min\{k,n\}$ with $j<N$, and return the first feasible
prefix. If none succeeds, no witness exists within budget $k$. One pass computes the
initial sum and prefix sums after sorting. The restriction $j<N$
ensures nonemptiness, including when every row is deletable.
\end{proof}
This case has no joins or shared supports. It illustrates why the failure
of AVG submodularity alone is not a hardness proof.
